% ===========================================================================
%  11 — Interface do usuario
% ===========================================================================
\chapter{Interface do usuário}
\label{cap:interface}

A interface é uma camada de apresentação estrita: consome exclusivamente a API
e não contém lógica de domínio nem acesso a banco (\adrr{ADR-06}). Este capítulo
documenta sua arquitetura de informação, os painéis analíticos e as decisões de
visualização --- que foram tomadas com validação, não por gosto.

\section{Arquitetura de informação}

\begin{figure}[H]
\centering
\fitwidth{%
\begin{tikzpicture}[x=1cm, y=1cm]

  \node[c4container, text width=28mm, minimum height=12mm] (nav) at (2.2,-3.6)
    {\textbf{Navegação lateral}\desc{seção + subpainel}};

  \node[c4component, text width=32mm, minimum height=11mm] (s1) at (7.4,0)
    {\textbf{Painéis}\desc{4 subpainéis analíticos}};
  \node[c4component, text width=32mm, minimum height=11mm] (s2) at (7.4,-2.4)
    {\textbf{Diagnóstico}\desc{JSON $\rightarrow$ resposta completa}};
  \node[c4component, text width=32mm, minimum height=11mm] (s3) at (7.4,-4.8)
    {\textbf{Conversa}\desc{sessão persistente, sugestões}};
  \node[c4component, text width=32mm, minimum height=11mm] (s4) at (7.4,-7.2)
    {\textbf{Documentos}\desc{lista + cadastro de PDF}};

  \node[c4component, text width=34mm, minimum height=9mm] (p1) at (13.6,1.3)
    {\textbf{Visão geral}\desc{distribuição, série temporal, priorização}};
  \node[c4component, text width=34mm, minimum height=9mm] (p2) at (13.6,-0.5)
    {\textbf{Severidade e física}\desc{RMS por rotação, zonas normativas}};
  \node[c4component, text width=34mm, minimum height=9mm] (p3) at (13.6,-2.3)
    {\textbf{Assinaturas}\desc{escores padronizados, poder discriminativo}};
  \node[c4component, text width=34mm, minimum height=9mm] (p4) at (13.6,-4.1)
    {\textbf{Qualidade do modelo}\desc{3 protocolos, importância, vazamento}};

  \node[externo, text width=26mm, minimum height=11mm] (st) at (2.2,-7.4)
    {\textbf{Estado do sistema}\desc{recolhido por padrão}};

  \draw[c4rel] (nav.east) -- (5.6,0);
  \draw[c4rel] (nav.east) -- (5.6,-2.4);
  \draw[c4rel] (nav.east) -- (5.6,-4.8);
  \draw[c4rel] (nav.east) -- (5.6,-7.2);
  \draw[c4rel] (s1.east) -- (11.6,1.3);
  \draw[c4rel] (s1.east) -- (11.6,-0.5);
  \draw[c4rel] (s1.east) -- (11.6,-2.3);
  \draw[c4rel] (s1.east) -- (11.6,-4.1);
  \draw[c4reldash] (nav.south) -- (st.north);

\end{tikzpicture}}
\caption[Arquitetura de informação da interface]{Arquitetura de informação. A
  navegação renderiza apenas a seção ativa, e cada painel requisita somente os
  dados de que precisa --- não há carga antecipada do conjunto completo.}
\label{fig:ia}
\end{figure}

\section{Seção de diagnóstico}

Recebe o JSON do evento --- com botão para carregar o exemplo do enunciado --- e
apresenta a resposta em quatro blocos:

\begin{enumerate}
  \item \textbf{Cabeçalho de métricas}: tipo de defeito, probabilidade,
    concordância dos vizinhos e selo de confiança, colorido conforme o nível.
  \item \textbf{Ocorrências históricas}: total registrado, frequência diária,
    período de observação e a distribuição temporal em gráfico de barras.
  \item \textbf{Instruções de correção} em Markdown, com as fontes citadas em
    seção expansível --- documento, seção e página.
  \item \textbf{Ou} o aviso de falha não documentada com a sugestão de registro,
    em destaque de atenção.
\end{enumerate}

Quando o diagnóstico aponta estado operacional, a interface informa que não há
ação prescritiva a tomar --- e não exibe seção de prescrição alguma, evitando a
sugestão implícita de que algo precisa ser feito.

\section{Painéis analíticos}
\label{sec:paineis}

Cada gráfico do painel entrou por responder a uma pergunta operacional
específica \emph{e} por ter passado no teste de não enganar quem o lê. A
fundamentação estatística está no Capítulo~\ref{cap:ml}; aqui registra-se a
tradução em produto.

\subsection{Priorização documental}
\label{sec:priorizacao}

\textbf{Pergunta.} O sistema só prescreve correção para falhas documentadas.
Quais lacunas custam mais caro?

\textbf{Resposta.} Uma matriz de bolhas que cruza três dimensões: frequência
(contagem por família) no eixo horizontal, severidade (máxima severidade
relativa entre regimes de rotação) no vertical, e cobertura documental na cor.
O canto superior direito em laranja é a fila de trabalho da engenharia de
manutenção.

\begin{tabularx}{\linewidth}{@{}L{34mm}R{26mm}X@{}}
\toprule
\headrow \thd{Família sem documento} & \thd{Ocorrências} & \thd{Situação}\\
\midrule
\cd{eccentric\_rotor} & \num{16497} & Maior lacuna do acervo --- segunda família
  mais frequente do histórico.\\
\zebra \cd{ventoinha} & \num{12299} & Sem procedimento cadastrado.\\
\cd{falta\_fase} & 800 & Sem procedimento, mas é a classe que o modelo
  reconhece melhor ($\text{F1} = 0{,}940$).\\
\bottomrule
\end{tabularx}

\subsection{Severidade e validação física}

O painel apresenta a velocidade eficaz \textbf{sempre estratificada por regime
de rotação}, com dois modos de leitura: absoluto, com as zonas normativas
sobrepostas e classe de máquina selecionável; e relativo à linha de base do
mesmo regime. O gráfico de validação física destaca o desbalanceamento contra as
demais famílias em cinza, evidenciando a dependência quadrática com a rotação
prevista pelo manual (Seção~\ref{sec:fisica}).

\subsection{Assinaturas}

Mapa de calor dos escores padronizados por família e por \textit{feature},
\textbf{ordenado pelo poder discriminativo} --- não pela ordem das colunas no
conjunto de dados. Indicadores fracos aparecem em cinza, com legenda explicando
o descompasso frente ao que os manuais previam
(Seção~\ref{sec:assinaturas}).

\subsection{Qualidade do modelo}

Apresenta os três protocolos de avaliação lado a lado, o F1 por classe, a
importância das \textit{features} com a temperatura destacada e o gráfico de
vazamento por \textit{feature}.

\begin{decisao}
Exibir os três protocolos é uma decisão de produto, não de estética. Um painel
que mostrasse apenas os \SI{89}{\percent} da divisão aleatória comunicaria uma
capacidade que o sistema não tem em produção --- e a primeira vez que um técnico
confiasse nele e errasse, a credibilidade da ferramenta inteira estaria perdida.
\end{decisao}

\section{Decisões de visualização}

As cores não foram escolhidas por preferência: cada conjunto passou por um
validador de paleta que verifica banda de luminosidade, piso de croma, separação
para deficiências de visão de cores e contraste contra a superfície, nos modos
claro e escuro. A diferença perceptual é medida em CIEDE2000~\cite{sharma2005ciede2000},
com piso de $\Delta E \geq 6$ sob simulação de deuteranopia, conforme a prática
recomendada para material acessível~\cite{okabe2008cud}.

\begin{table}[H]
\caption{Codificações visuais e sua justificativa}
\label{tab:visualizacao}
\begin{tabularx}{\linewidth}{@{}L{32mm}L{38mm}X@{}}
\toprule
\headrow \thd{Codificação} & \thd{Escolha} & \thd{Justificativa}\\
\midrule
Documentada $\times$ sem documento & Azul \cd{\#2a78d6} / laranja
  \cd{\#eb6834} &
  A escolha intuitiva verde/vermelho \textbf{falhou} o teste: $\Delta E = 4{,}1$
  sob deuteranopia, abaixo do piso de 6. Azul/laranja atinge
  $\Delta E = 24{,}7$. A legenda usa ainda ícone como codificação
  redundante.\\
\zebra Regime de rotação & Rampa ordinal de um único tom (azul, claro para
  escuro) &
  Rotação é grandeza \textbf{ordenada}, não identidade. Matizes categóricos
  sugeririam categorias sem ordem, o que é uma afirmação
  falsa~\cite{munzner2014visualization}.\\
Escore padronizado & Divergente azul $\leftrightarrow$ vermelho, cinza no meio &
  O escore expressa \textbf{polaridade} --- acima ou abaixo da média global.
  Uma escala arco-íris destruiria essa leitura.\\
\zebra Famílias no gráfico de física & Cinza mais um destaque &
  Doze séries excedem qualquer paleta categórica segura. Destaque sobre
  referência cinza comunica melhor do que doze cores que ninguém
  distingue.\\
Severidade & Percentis (50, 95, 99), nunca média &
  As distribuições têm cauda longa; a média seria puxada por transientes de
  partida e daria uma impressão falsa de severidade típica.\\
\bottomrule
\end{tabularx}
\end{table}

\begin{nota}
O registro do \emph{erro evitado} é a parte mais útil desta tabela. A
codificação verde/vermelho para \enquote{documentado / não documentado} é a
primeira escolha de qualquer pessoa, é usada em incontáveis painéis e é
inacessível para aproximadamente \SI{8}{\percent} dos homens. O teste custou
minutos; a correção, uma linha de código.
\end{nota}

\section{Estado do sistema}

A barra lateral apresenta, em bloco recolhido por padrão, o estado corrente:
modelo de linguagem e de \textit{embedding} configurados, presença dos artefatos
de aprendizado de máquina e número de famílias documentadas. Se a API estiver
inacessível, a interface interrompe a renderização com mensagem acionável
contendo o endereço configurado --- em vez de exibir painéis vazios que parecem
indicar ausência de dados.
